\section{The cases \texorpdfstring{$9\le n\le14$}{9<=n<=14}}
\label{sec:small-uniform-en}

This section separates the two larger cases, which are decided at the level of
global integer data, from the four smaller cases that require an exact finite
verification.  No historical search tree or saved orbit list is used.

For $k\ge3$, let $\ell_k$ and $c_k$ denote the numbers of maximal $k$-point
lines and circles, and put $m_k=\ell_k+c_k$.
Every pair belongs to one maximal
line, while every triple belongs either to its maximal line or to its unique
circumcircle.  Therefore
\begin{align}
 \ell_2+\sum_{k\ge3}\binom{k}{2}\ell_k&=\binom n2,
 \label{eq:small-line-pairs-en}\\
 \sum_{k\ge3}\binom{k}{3}m_k&=\binom n3,
 \qquad \sum_{k\ge3}c_k=c(P).
 \label{eq:small-triple-partition-en}
\end{align}
We call $(\ell_3,\ldots,\ell_K;c_3,\ldots,c_K)$ the \emph{global count
vector}; it records only the number of maximal lines and circles of each size,
not the subsets of points on them.  Lemma~\ref{lem:largest-en} excludes maximum blocks of size at
least $6,7,7,7,7,8$ for $n=9,10,11,12,13,14$, respectively.  Thus $K$ is at
most $5,6,6,6,6,7$ in the branches considered below.

The enumeration of these count vectors uses only necessary conditions.  Since $K<n-1$,
every one-point deletion is admissible, and summing its circle count gives
\begin{equation}\label{eq:small-one-delete-en}
 n c(P)-3c_3\ge n c(n-1).
\end{equation}
Inversion followed by Zhang's prescribed-point ordinary-line theorem
\cite[Theorem~4.1]{zhang-en} gives
\begin{equation}\label{eq:small-z-circle-en}
 3c_3\ge n\left\lceil\frac{n-1}{6}\right\rceil.
\end{equation}
We also use the ordinary-line, Melchior, Hirzebruch--Bojanowski, and
Shnurnikov inequalities for the line part of the count vector, and their sums
over subsets that are forced to be noncollinear
\cite[Eq.~(9)]{hirzebruch-en}
\cite[Theorem~2.1]{pokora-en}\cite[Theorem~3]{shnurnikov-en}.
For a subset contained in a maximal line no circle lower bound is imposed; for
a subset contained in a maximal circle only that defining circle is imposed.
This is the correction needed when deletion averages are used at small order.

We shall also use the following elementary consequence of block intersections.
For $N\ge1$ and $S=aN+b$, $0\le b<N$, set
\[
 \Phi_N(S)=N\binom a2+ab.
\]
This is the minimum of $\sum_i\binom{x_i}{2}$ over nonnegative integral
$x_i$ with $\sum_i x_i=S$.  If a selected family consists of $L$ maximal
lines and $G$ maximal circles and $I_1,I_2$ are its total point and point-pair
incidences, then
\begin{align}
 \Phi_n(I_1)&\le \binom L2+2LG+2\binom G2,
 \label{eq:small-block-point-en}\\
 \Phi_{\binom n2}(I_2)&\le LG+\binom G2.
 \label{eq:small-block-pair-en}
\end{align}
Indeed, two lines meet in at most one point and any other two distinct blocks
meet in at most two points.  Applying these inequalities to the largest
members of each subfamily covers every possible choice.  It does not decide
which labelled points lie on the selected lines and circles.

\subsection{The augmented-line inequality}

Fix $x\in P$, invert about $x$, and dualize the remaining $n-1$ points.  Let
$t_j(x)$ be the number of multiplicity-$j$ vertices of the resulting
arrangement and put
\[
 \delta_x=t_2(x)-3-\sum_{j\ge4}(j-3)t_j(x).
\]
Write $d_k^L(x)$ for the number of maximal $k$-point lines of $P$ through
$x$.

\begin{lemma}\label{lem:augmented-line-en}
For every $x\in P$,
\begin{equation}\label{eq:augmented-line-point-en}
 \sum_{k\ge3}k d_k^L(x)\le n-1+\delta_x.
\end{equation}
Consequently,
\begin{equation}\label{eq:augmented-line-summed-en}
 \sum_{k\ge3}k^2\ell_k\le n(n-1)+D,
 \qquad
 D=3m_3-3n-\sum_{k\ge5}k(k-4)m_k.
\end{equation}
\end{lemma}

\begin{proof}
Let $\mathcal A_x$ be the arrangement of $n-1$ dual lines, and adjoin the
projective line dual to the inversion centre.  A $k$-point original line
through $x$ corresponds to an old vertex of multiplicity $k-1$ on the added
line.  Suppose the selected old vertices have multiplicities $i_1,\ldots,i_s$.
The added line has
$(n-1)-\sum_j i_j$ new ordinary vertices.  Raising the multiplicity of each
selected old vertex by one decreases the Melchior defect by one, including
the cases of multiplicity two and three.  Hence the defect of the enlarged
arrangement is
\[
 \delta_x+(n-1)-\sum_{j=1}^s(i_j+1)
 =\delta_x+(n-1)-\sum_{k\ge3}k d_k^L(x).
\]
The enlarged arrangement is essential, so Melchior's inequality proves
\eqref{eq:augmented-line-point-en}.  Summing over $x$, and using
$\sum_xd_k^L(x)=k\ell_k$ and the local incidence identities, gives
\eqref{eq:augmented-line-summed-en}.
\end{proof}

\subsection{Global exclusion for thirteen and fourteen points}

For $n=14$, the exact integer solutions of the preceding necessary
conditions with $c(P)<F(14)$ give 8981 count vectors.  Inequality
\eqref{eq:augmented-line-summed-en} excludes 8980.  The remaining vector is
\[
 (\ell_3,\ldots,\ell_7;c_3,\ldots,c_7)
 =(20,0,0,0,0;23,44,0,2,3).
\]
Its three seven-point circles have pairwise intersections of size at most two,
so their union has at least $3\cdot7-3\cdot2=15$ points, a contradiction.
Thus $c(14)=F(14)$.

For $n=13$, the same exact enumeration gives 3962 count vectors below $F(13)$.
The augmented-line inequality leaves 137, and
\eqref{eq:small-block-point-en}--\eqref{eq:small-block-pair-en} leave 73.  We
now derive two inequalities that exclude these remaining vectors.

At any point $x$, the local arrangement has twelve lines and, in the present
branch, no vertex of multiplicity greater than five.  Write its multiplicity
vector as $(t_2,t_3,t_4,t_5)$ and its defect as $\delta$.  Then
\begin{equation}\label{eq:n13-local-two-ineq-en}
 -2\delta+t_4+3t_5\le6,
 \qquad -\delta+t_5\le1.
\end{equation}
For the first inequality, the pair identity and Bojanowski's inequality give
\begin{align*}
 t_2+3t_3+6t_4+10t_5&=66,\\
 4t_2+3t_3&\ge48+5t_5,
\end{align*}
and hence
\[
 H:=3t_2-6t_4-15t_5+18\ge0.
\]
Put $E=2t_2-3t_4-7t_5$.  Since
$H+3\delta=3(E+3)$, one has $E\ge-3$.  Moreover, the pair identity gives
$E\equiv0\pmod3$.  Equality $E=-3$ would force $H=\delta=0$.
Cuntz's complete list of simplicial twelve-line arrangements, after excluding
a six-fold vertex, leaves the types $(8,10,3,1)$ and $(9,7,6,0)$; both have
$E=0$ \cite[Theorem~1.1 and the invariant tables]{cuntz-en}.  Therefore
$E\ge0$, which is the first inequality in \eqref{eq:n13-local-two-ineq-en}.

For the second inequality, four five-fold vertices would dualize to four
five-point lines whose union has at least $20-\binom42=14$ points; hence
$t_5\le3$.  If $\delta\ge2$ the claim follows.  If $\delta=0$, Cuntz's two
types above have $t_5\le1$.  The only remaining possible violation is
$(\delta,t_5)=(1,3)$.  The pair and defect identities then force
$(t_2,t_3,t_4,t_5)=(12,4,2,3)$.  But three five-point lines and one
four-point line have a union of at least $5+5+5+4-\binom42=13$ points, again
impossible on twelve points.  This proves the second inequality.

Summing \eqref{eq:n13-local-two-ineq-en} over the thirteen centres yields
\begin{equation}\label{eq:n13-global-two-ineq-en}
 -2D+5m_5+18m_6\le78,
 \qquad -D+6m_6\le13.
\end{equation}
Of the 73 count vectors, 71 violate the first inequality and the other two
violate the second.  This is an exact substitution into
\eqref{eq:n13-global-two-ineq-en}; no subsets of labelled points are selected.  Hence
$c(13)=F(13)$.

\subsection{Exact verification for nine through twelve points}

For $9\le n\le12$, the same global constraints and
\eqref{eq:augmented-line-summed-en} leave a small but nonempty boundary.  We
describe the common verification so that every transition can be reproduced.
For a marked point $x$, let $a_k(x)$ and $u_k(x)$ count maximal $k$-point
lines and circles through $x$.  Then
\[
 t_{k-1}(x)=a_k(x)+u_k(x),
 \qquad
 \sum_{j\ge2}\binom j2t_j(x)=\binom{n-1}{2}.
\]
These local multiplicity vectors obey the published arrangement inequalities used above, the
block-union inequalities, and \eqref{eq:augmented-line-point-en}.  A first
integer model records only how many points have each vector and imposes
\[
 \sum_xa_k(x)=k\ell_k,\qquad
 \sum_xu_k(x)=kc_k,\qquad
 \sum_x\delta_x=D.
\]

The next two stages still do not assign labelled point subsets.  For a block category
$\alpha$, let $d_\alpha(x)$ be its point degree and let
$\lambda_\alpha(e)$ be its degree at a point pair $e$.  If $k_\alpha$ is the
block size and $N_\alpha$ its multiplicity, then
\begin{align}
 \sum_xd_\alpha(x)&=k_\alpha N_\alpha,&
 \sum_e\lambda_\alpha(e)&=\binom{k_\alpha}{2}N_\alpha,
 \label{eq:small-category-totals-en}\\
 \sum_\alpha(k_\alpha-2)\lambda_\alpha(e)&=n-2.
 \label{eq:small-pair-profile2-en}
\end{align}
The last identity says that the portions outside $e$ of all maximal blocks
through $e$ partition $P\setminus e$.  The point and point-pair data are linked by
\[
 \sum_{y\ne x}\lambda_\alpha(\{x,y\})=(k_\alpha-1)d_\alpha(x).
\]
The subsequent intersection count uses,
for two block categories, the exact identities
\[
 N_2=M_2,\qquad N_1=M_1-2M_2,\qquad N_0=Q-M_1+M_2,
\]
where $M_1$ is the sum of intersection sizes, $M_2$ is the sum of their second
binomial coefficients, and $N_s$ counts pairs of
distinct blocks meeting in $s$ points.  It also counts triples disjoint from
a distinguished block.  All coefficients are binomial counts, and no
floating-point relaxation is used.

Finally, inversion at $x$ determines the complete count vector on $n-1$ image
points: blocks through $x$ become image lines with one fewer point, whereas
blocks avoiding $x$ become image circles.  A local vector is retained only if this
child vector survives the next level, after which the local multiplicity data
at the $n$ points are balanced again.  The recursion stops at $n=9$.  There the model
selects maximal line and circle subsets of the nine labelled points, requires
every triple to lie in exactly one selected block, requires a point pair to
lie in at most one selected line, and imposes the already proved lower bounds
on every admissible proper subset.  Infeasibility of this more permissive
abstract model implies geometric infeasibility.

The exact counts in the final replay are:
\begin{center}
\small
\begin{tabular}{c|rrrrrrr}
\toprule
$n$ & initial & augmented sum & local vectors & point--pair & intersections & block pairs & terminal\\
\midrule
9  &45   &36  &23  &-- &-- &-- &0\\
10 &180  &150 &55  &26 &21 &12 &1\\
11 &416  &160 &70  &39 &28 &-- &0\\
12 &1354 &211 &169 &55 &44 &-- &0\\
\bottomrule
\end{tabular}
\end{center}
Here ``initial'' is the number left from 129, 551, 1215, and 4322 raw
count vectors, respectively.  The terminal zeroes for $n=11,12$ come from the
same recursion.  It visits, respectively,
\[
 (28,72,222)\quad\text{and}\quad(44,322,661,622)
\]
count vectors at orders $(11,10,9)$ and $(12,11,10,9)$.  Every excluded status
is an exact \texttt{INFEASIBLE} status; there is no \texttt{UNKNOWN}.

For $n=10$ the recursion leaves the single count vector
\[
 (\ell_3,\ell_4,\ell_5,\ell_6;
   c_3,c_4,c_5,c_6;\ell_2)
 =(10,0,0,0;10,20,2,0;15).
\]
The four local multiplicity vectors retained by the recursive certificate all have
$u_5(x)=1$ and $a_3(x)\le3$.  Since the two five-point circles have ten point
incidences, they therefore partition the ten points into sets $A$ and $B$.
Moreover, the ten three-point lines have thirty point incidences, so
$a_3(x)=3$ at every point; the corresponding retained vector also has
$u_4(x)=8$.  Every four-point circle meets each of $A$ and $B$ in two points.

Associate such a circle with its chord in $A$ and its chord in $B$.  The two
chords meet on the fixed radical axis of the two five-point circles.  No three
chords determined by five points on a circle can pass through the same point
off that circle: three pairwise disjoint chords would require six endpoints.
Thus each set of chords meeting at one radical-axis point has size at most two
on either side.  The twenty four-point circles attain the resulting maximum,
so the chord graph is five disjoint copies of $K_{2,2}$.  Up to the action of
$S_5\times S_5$, the six perfect matchings of the Petersen graph give two
double orbits.  One admits no compatible ten-line design; the other has two
complementary designs forming one orbit.

For a representative of that orbit, the ten line incidences give a complete
projective parametrization with nonzero parameters $a,b,d$.  Five selected
circle-intersection minors are
\begin{gather*}
 ad+b^2,\quad -ab+ad-b,\quad ab-ad-a-1,\\
 -ad+bd-b-d,\quad b(ab+1).
\end{gather*}
They must all vanish.  Adding the second and third relations gives
$a+b+1=0$, while the last relation and $b\ne0$ give $ab+1=0$.  Hence
$b^2+b-1=0$.  The first relation gives
$d=b^2/(b+1)$; here $b+1=-a\ne0$.  Substitution in the fourth relation then
gives, after removal of nonzero factors,
\[
 (b-1)(2b+1)=0,
\]
which is relatively prime to $b^2+b-1$.  The discarded parameter branch makes
two labelled points coincide, and every factor divided out in forming the
five primitive minors is one of $a,d,b$ or $ad$, all nonzero.  Thus the last
$n=10$ count vector is impossible.

Combining these verifications with the construction in \eqref{eq:F-en} gives
\[
 c(n)=F(n)\qquad(9\le n\le14).
\]

\begin{remark}[Scope of the finite verification]
The only labelled search in this section occurs in the nine-point terminal
model.  Orders ten through twelve use count vectors and recursive consistency
of local multiplicities, while orders thirteen and fourteen use only global inequalities.
No condition $\sum_jt_j\ge2r-4$ is assumed for an arbitrary arrangement.
For an arrangement of $q$ lines, Elliott's Theorem~3 is used only when its
hypotheses $q\ge10$ and no $(q-1)$-fold vertex are explicitly satisfied.
\end{remark}
